\documentclass[11pt]{article}
\usepackage[localtheorems]{raeez-math-template}

\newtheorem{theorem}{Theorem}[section]
\newtheorem{proposition}[theorem]{Proposition}
\newtheorem{lemma}[theorem]{Lemma}
\newtheorem{corollary}[theorem]{Corollary}
\theoremstyle{definition}
\newtheorem{definition}[theorem]{Definition}
\newtheorem{remark}[theorem]{Remark}
\newtheorem{conjecture}[theorem]{Conjecture}

\newcommand{\C}{\mathbb{C}}
\newcommand{\Z}{\mathbb{Z}}
\newcommand{\Q}{\mathbb{Q}}
\newcommand{\F}{\mathcal{F}}
\renewcommand{\gg}{\mathfrak{g}}
\newcommand{\hh}{\mathfrak{h}}
\newcommand{\Lmuk}{\Lambda_{\mathrm{Muk}}}
\newcommand{\YgK}{Y_\hbar(\gg_{K3})}
\newcommand{\Yplus}{Y_\hbar^{+}(\gg_{K3})}
\newcommand{\Yminus}{Y_\hbar^{-}(\gg_{K3})}
\newcommand{\Hilb}{\mathrm{Hilb}}
\newcommand{\OPE}{\mathrm{OPE}}
\newcommand{\RTT}{\mathrm{RTT}}
\newcommand{\Rep}{\mathrm{Rep}}
\newcommand{\End}{\mathrm{End}}
\newcommand{\CoHA}{\mathrm{CoHA}}
\newcommand{\Sym}{\mathrm{Sym}}
\newcommand{\Obs}{\mathrm{Obs}}
\newcommand{\Tr}{\mathrm{Tr}}
\newcommand{\Coh}{\mathrm{Coh}}
\newcommand{\Aut}{\mathrm{Aut}}

\title{The K3 Abelian Yangian:\\
Mukai-Heisenberg Presentation, Local Stable Envelopes,\\
and the Holomorphic Chern-Simons Boundary}
\author{Raeez Lorgat}
\date{2026-04-22}

\begin{document}
\maketitle

\begin{abstract}
Let $S$ be a K3 surface and
\[
  \Lmuk=\widetilde H(S,\Z)=H^0(S,\Z)\oplus H^2(S,\Z)\oplus H^4(S,\Z)
  \simeq U^3\oplus E_8(-1)^2
\]
the Mukai lattice with pairing of signature $(4,20)$. The rank $24$
comes from $1+22+1$; the middle cohomology alone has rank $22$ and
signature $(3,19)$. The abelian K3 Yangian $\YgK$ is the rational
Heisenberg Yangian attached to the rank-$24$ Mukai current algebra. Its
strict algebraic core is the OPE
\[
  J_i(z)J_j(w)\sim \frac{\hbar\,\omega^{ij}}{(z-w)^2},
\]
the positive diagonal RTT half, and the Heisenberg double pairing between
positive and negative modes. The diagonal spectral data are the scalar
family
\[
  r_i(u)=\frac{u-h_i}{u+h_i}, \qquad
  g_{K3}(u)=\prod_{i=1}^{24}\frac{u-h_i}{u+h_i},
  \qquad \sum_i h_i=0,
\]
so the structure function has bidegree $(24,24)$ and satisfies
$g_{K3}(u)g_{K3}(-u)=1$. Local ADE/Kummer charts supply the
Schiffmann-Vasserot quiver-CoHA positive half and the Maulik-Okounkov
stable-envelope $R$-matrix. Generic K3 moduli lack the rank-$2$ torus
needed for a global MO stable envelope, and the compact $K3\times E$
Hall-Drinfeld/BKM branch is a separate construction. The hCS boundary
interpretation is stated as a conjectural avatar of the same
Mukai-Heisenberg algebra.
\end{abstract}

\tableofcontents

\section{Mukai-Heisenberg input}
\label{sec:mukai-heisenberg-input}

\subsection{The rank is Mukai rank}

The K3 lattice $H^2(S,\Z)$ has rank $22$ and signature $(3,19)$. The
Mukai lattice
\[
  \Lmuk=H^0(S,\Z)\oplus H^2(S,\Z)\oplus H^4(S,\Z)
\]
adds the two degree $0$ and degree $4$ directions, giving rank $24$ and
signature $(4,20)$. The $24$ currents of the K3 abelian Yangian are
indexed by the full Mukai lattice rather than by $H^2(S,\Z)$ alone.

Fix a diagonal real basis $\{e_i\}_{i=1}^{24}$ for
$\Lmuk\otimes\C$ in which
\[
  \omega(e_i,e_j)=\omega^{ij}=\varepsilon_i\delta_{ij},
  \qquad
  \varepsilon_i=
  \begin{cases}
    +1,&1\leq i\leq4,\\
    -1,&5\leq i\leq24.
  \end{cases}
\]
The diagonal basis is a computational normal form; the signature and the
Heisenberg algebra are basis-invariant.

\begin{definition}[Abelian K3 Yangian]
\label{def:abelian-k3-yangian}
The abelian K3 Yangian $\YgK$ is the rational RTT/Drinfeld-current
Heisenberg Yangian generated by modes
\[
  \{J_{i,n}\}_{1\leq i\leq24,\ n\in\Z}
\]
and a central element $\mathbf c$, with Heisenberg relation
\begin{equation}
\label{eq:heis-commutator}
  [J_{i,m},J_{j,n}]
  =
  \hbar\,\omega^{ij}\,m\,\delta_{m+n,0}\,\mathbf c.
\end{equation}
The central charge is normalised by $\mathbf c=1$. The diagonal spectral
parameters $h_1,\ldots,h_{24}\in\C$ enter the rational $R$-matrix below
and satisfy the trace-zero CY$_2$ normalisation
\begin{equation}
\label{eq:cy2-trace-zero}
  \sum_{i=1}^{24}h_i=0.
\end{equation}
\end{definition}

The current
\[
  J_i(z)=\sum_{n\in\Z}J_{i,n}z^{-n-1}
\]
has OPE
\begin{equation}
\label{eq:k3-heisenberg-ope}
  J_i(z)J_j(w)\sim \frac{\hbar\,\omega^{ij}}{(z-w)^2}.
\end{equation}
Equivalently, in Lie conformal notation,
\[
  [J_i{}_\lambda J_j]=\hbar\,\omega^{ij}\lambda.
\]
This is the rank-$24$ Mukai-Heisenberg vertex algebra
$V_{\Lmuk}$ restricted to its Cartan current subalgebra.

\subsection{Miura series and diagonal spectral data}

Use the normalised Miura series
\begin{equation}
\label{eq:miura-normalised}
  \mathcal T_{K3}(u)
  =
  \prod_{i=1}^{24}\left(1-\frac{\phi_i}{u}\right)
  =
  1+\sum_{s=1}^{24}(-1)^s\Psi_s u^{-s},
\end{equation}
where $\phi_i=J_i+\hbar\partial_z$ is the $i$-th free field. The
unnormalised polynomial transfer matrix is
$T_{K3}(u)=u^{24}\mathcal T_{K3}(u)$. The coefficients
$\Psi_s=e_s(\phi_1,\ldots,\phi_{24})$ are elementary symmetric Miura
fields. The truncation $\Psi_s=0$ for $s>24$ is the elementary identity
$e_s(x_1,\ldots,x_{24})=0$.

The structure function is
\begin{equation}
\label{eq:structure-function-k3}
  g_{K3}(u)=\prod_{i=1}^{24}\frac{u-h_i}{u+h_i}.
\end{equation}
It has bidegree $(24,24)$, normalisation $g_{K3}(\infty)=1$, and
unitarity $g_{K3}(u)g_{K3}(-u)=1$. The coefficient of $u^{-1}$ in
$\log g_{K3}(u)$ is $-2\sum_i h_i$, hence \eqref{eq:cy2-trace-zero}
removes the trace term in the diagonal Cartan.

On the Mukai weight basis $V=\C^{24}$ define the diagonal scalar
$R$-family by
\begin{equation}
\label{eq:r-matrix-diagonal}
  R_{K3}(u)(e_i\otimes e_j)
  =
  \frac{u-h_i}{u+h_i}\,e_i\otimes e_j.
\end{equation}
Thus $R_{K3}(u)\in\End(V\otimes V)[[u^{-1}]]$ is diagonal with each
$r_i(u)$ repeated along the second tensor factor. The Yang-Baxter
equation follows entrywise from commutativity of scalar diagonal entries:
\[
  R_{12}(u-v)R_{13}(u)R_{23}(v)
  =
  R_{23}(v)R_{13}(u)R_{12}(u-v).
\]
At large $u$,
\begin{equation}
\label{eq:classical-r-mukai-note}
  R_{K3}(u)=I-\frac{2}{u}\operatorname{diag}(h_1,\ldots,h_{24})+O(u^{-2}).
\end{equation}

\section{RTT half and Heisenberg double}
\label{sec:rtt-half}

\begin{proposition}[Positive diagonal RTT half]
\label{prop:rtt-positive-half}
Let
\[
  L^+(u)=\operatorname{diag}(L^+_1(u),\ldots,L^+_{24}(u)),
  \qquad
  L^+_i(u)=1+\sum_{r\geq0}\hbar J_{i,r}u^{-r-1}.
\]
The diagonal RTT equation
\begin{equation}
\label{eq:rtt-relation}
  R_{K3,12}(u-v)L^+_1(u)L^+_2(v)
  =
  L^+_2(v)L^+_1(u)R_{K3,12}(u-v)
\end{equation}
is equivalent to commutativity of the nonnegative modes
\[
  [J_{i,r},J_{j,s}]=0,\qquad r,s\geq0.
\]
It presents the positive current half. The full relation
\eqref{eq:heis-commutator} is obtained by adjoining the negative half and
the Mukai Heisenberg pairing.
\end{proposition}

\begin{proof}
All matrices in \eqref{eq:rtt-relation} are diagonal. Comparing the
$(i,j)$ component gives
\[
  r_i(u-v)L^+_i(u)L^+_j(v)
  =
  L^+_j(v)L^+_i(u)r_i(u-v).
\]
The scalar $r_i(u-v)$ cancels, so the condition is
$L^+_i(u)L^+_j(v)=L^+_j(v)L^+_i(u)$ for all $i,j$. Coefficient extraction
gives $[J_{i,r},J_{j,s}]=0$ for all $r,s\geq0$. Conversely these
commutators make the diagonal RTT equation an identity.

The central extension in \eqref{eq:heis-commutator} pairs positive and
negative modes:
\[
  [J_{i,r},J_{j,-s}]
  =
  \hbar\,\omega^{ij}r\,\delta_{r,s},
  \qquad r,s>0.
\]
This is the Heisenberg double of the two commutative current halves.
\end{proof}

\begin{remark}[Rank two test]
\label{rem:rank-2-rtt}
For the rank-$2$ diagonal form
$\operatorname{diag}(+1,-1)$ the RTT equation states that
$L^+_1(u)$ and $L^+_2(v)$ commute. The Heisenberg double adds
$[J_{1,r},J_{1,-r}]=\hbar r$ and
$[J_{2,r},J_{2,-r}]=-\hbar r$. The parameters $h_1,h_2$ act through
the scalar $R$-entries and through the shifted coproduct, while the
central extension is controlled by the Mukai form.
\end{remark}

\begin{remark}[Yangian terminology]
\label{rem:yangian-terminology}
The word \emph{Yangian} is used in the rational RTT/Drinfeld-current
sense: a spectral-parameter Hopf algebra with rational $R$-matrix,
Yang-Baxter equation, triangular PBW decomposition, and Drinfeld double
pairing. The simple-Lie Chevalley-Serre package belongs to ADE
enhancement loci. At the generic Mukai-Heisenberg point the Lie bracket
is abelian modulo the central line and the rational spectral structure is
carried by $R_{K3}(u)$ and $\Delta_z$.
\end{remark}

\subsection{Current presentation}

\begin{proposition}[Mukai current presentation]
\label{prop:mukai-current-presentation}
The current algebra presentation of $\YgK$ is generated by
$J_i(z)$, $1\leq i\leq24$, with lambda brackets
\[
  [J_i{}_\lambda J_j]=\hbar\,\omega^{ij}\lambda.
\]
Its triangular decomposition is
\[
  \YgK\simeq \YgK^-\otimes\YgK^0\otimes\YgK^+,
  \qquad
  \YgK^\pm=\C[J_{i,\mp n}:1\leq i\leq24,\ n>0],
  \qquad
  \YgK^0=\C[J_{1,0},\ldots,J_{24,0}].
\]
\end{proposition}

\begin{proof}
The lambda bracket is equivalent to \eqref{eq:heis-commutator}. PBW for
the Heisenberg algebra orders negative modes, zero modes, and positive
modes. The parameters $h_i$ enter the spectral and coproduct structure,
while the PBW basis is controlled only by the nondegenerate Mukai
pairing.
\end{proof}

\begin{remark}[Lie level and vertex level]
\label{rem:abelian-at-lie}
The derived subalgebra of
\[
  \bigoplus_{n\in\Z}\Lmuk\otimes t^n\oplus\C\mathbf c
\]
is the central line $\C\mathbf c$. Nonabelian brackets enter after
Frenkel-Kac vertex closure inside the lattice VOA. The BKM superalgebra
$\gg_{\Delta_5}$ belongs to that vertex-operator branch; the abelian K3
Yangian $\YgK$ belongs to the Mukai-Heisenberg branch.
\end{remark}

\section{Spectral Hopf structure}
\label{sec:hopf-structure}

\begin{proposition}[Shifted coproduct]
\label{prop:coproduct-drinfeld}
For each $z\in\C$, the formula
\begin{equation}
\label{eq:coproduct-drinfeld}
  \Delta_z\bigl(\mathcal T_{K3}(u)\bigr)
  =
  \mathcal T_{K3}^{L}(u)\mathcal T_{K3}^{R}(u-z)
\end{equation}
defines an algebra homomorphism
$\Delta_z:\YgK\to\YgK\otimes\YgK$ on the Miura generators. The spin-one
current is primitive,
\[
  \Delta_z(\Psi_1)=\Psi_1^L+\Psi_1^R,
\]
and the spin-two coefficient is
\[
  \Delta_z(\Psi_2)
  =
  \Psi_2^L+\Psi_2^R+\Psi_1^L\Psi_1^R+z\Psi_1^R.
\]
Higher coefficients are obtained by expanding $(u-z)^{-b}$ in powers of
$u^{-1}$.
\end{proposition}

\begin{proof}
Equation \eqref{eq:coproduct-drinfeld} multiplies two normalised
series with constant term $1$, so coefficient extraction gives a
well-defined algebra map on the polynomial Miura generators. The identity
\[
  (u-z)^{-b}
  =
  \sum_{p\geq0}\binom{b+p-1}{p}z^p u^{-b-p}
\]
gives the displayed spin-one and spin-two formulas. Since the Miura
fields are free Heisenberg fields acting on separate tensor factors,
normal ordering is preserved by tensor multiplication.
\end{proof}

\begin{proposition}[Coassociativity, counit, antipode]
\label{prop:hopf-axioms}
The coproducts satisfy
\[
  (\Delta_w\otimes\mathrm{id})\circ\Delta_{z+w}
  =
  (\mathrm{id}\otimes\Delta_z)\circ\Delta_w.
\]
The counit is determined by
$\epsilon(\mathcal T_{K3}(u))=1$ and $\epsilon(J_{i,n})=0$. The antipode
on the completed Miura algebra is determined by
\[
  S\bigl(\mathcal T_{K3}(u)\bigr)=\mathcal T_{K3}(u)^{-1},
  \qquad S(J_{i,n})=-J_{i,n}
\]
on primitive current modes.
\end{proposition}

\begin{proof}
Both sides of coassociativity send $\mathcal T_{K3}(u)$ to
\[
  \mathcal T^{(1)}(u)\mathcal T^{(2)}(u-w)\mathcal T^{(3)}(u-w-z).
\]
The counit identities follow because the normalised Miura series has
constant term $1$. The antipode is the inverse for multiplication in the
completed algebra of formal series with constant term $1$; on primitive
spin-one currents this is the sign map.
\end{proof}

\begin{remark}[Quasitriangularity]
\label{rem:quasi-triangular}
The diagonal $R$-matrix intertwines $\Delta_z$ with the opposite shifted
coproduct entrywise:
\[
  R_{K3}(u)\Delta_z^{\mathrm{op}}(x)=\Delta_z(x)R_{K3}(u).
\]
The off-diagonal nonabelian braiding is supplied by ADE enhancement or by
Frenkel-Kac vertex closure rather than by the diagonal Mukai-Heisenberg
current algebra itself.
\end{remark}

\section{Frenkel-Kac vertex closure}
\label{sec:fk-closure}

For $\alpha\in\Lmuk$ define the lattice vertex operator
\begin{equation}
\label{eq:vertex-operator-fk}
  V_\alpha(z)
  =
  \exp\!\left(\sum_{n\geq1}\frac{\alpha_{-n}}{n}z^n\right)
  \exp\!\left(-\sum_{n\geq1}\frac{\alpha_n}{n}z^{-n}\right)
  e^\alpha z^{\alpha_0}.
\end{equation}
The Frenkel-Kac bimultiplicative cocycle
$\epsilon_{\mathrm{FK}}:\Lmuk\times\Lmuk\to\{\pm1\}$ gives
\[
  V_\alpha(z)V_\beta(w)
  =
  \epsilon_{\mathrm{FK}}(\alpha,\beta)
  (z-w)^{\omega(\alpha,\beta)}
  :V_\alpha(z)V_\beta(w):
  +\cdots.
\]
For $\omega(\alpha,\beta)=-1$ the residue bracket is
\[
  [e_\alpha,e_\beta]
  =
  \epsilon_{\mathrm{FK}}(\alpha,\beta)e_{\alpha+\beta}.
\]
This is the vertex-operator source of the nonabelian BKM brackets. The
local anchor is
\texttt{chapters/examples/k3\_yangian\_chapter.tex},
Theorem \texttt{thm:k3-yangian-frenkel-kac-closure}.

\section{Local CoHA and stable envelopes}
\label{sec:local-coha-stable}

\subsection{Quiver CoHA positive half}

Let $\Gamma\subset SL_2(\C)$ be an ADE finite subgroup and let
$Q_\Gamma$ be the McKay quiver. The Schiffmann-Vasserot theorem
\texttt{arXiv:1705.07488} concerns the cohomological Hall algebra of the
preprojective quiver stack attached to $Q_\Gamma$ and its action on
Nakajima quiver varieties. It supplies the ADE chart input for K3
orbifold points. It is a quiver theorem, so the global compact K3 surface
requires separate gluing data.

\begin{proposition}[Local positive half]
\label{prop:local-positive-half}
On an ADE/Kummer chart of K3 moduli, the quiver CoHA of the
preprojective algebra of $Q_\Gamma$ gives the positive half of the local
Yangian acting on the corresponding Nakajima variety. After embedding the
ADE root lattice into $\Lmuk$, the Cartan part is the corresponding
Mukai sublattice of the rank-$24$ Heisenberg current algebra.
\end{proposition}

\begin{proof}
Schiffmann-Vasserot construct generators for the quiver CoHA and its
action on Nakajima quiver varieties. The McKay correspondence identifies
the exceptional curves of the ADE surface resolution with the simple
roots of $Q_\Gamma$. The inclusion of this root lattice into the Mukai
lattice places the local positive half inside the Mukai-Heisenberg
positive current half. The passage from this chartwise statement to a
compact K3 CoHA demands gluing over the period domain and a compact
orientation datum.
\end{proof}

The chartwise shuffle product has the form
\begin{equation}
\label{eq:shuffle-product}
  F_n(z;\alpha)G_m(w;\beta)
  =
  \prod_{a,b}g_{\Gamma}(z_a-w_b;\alpha_a,\beta_b)
  \operatorname{sh}(F\otimes G).
\end{equation}
For the abelian Mukai model, $g_\Gamma$ is replaced by the diagonal
$g_{K3}$ of \eqref{eq:structure-function-k3} after summing over the
twenty-four Mukai directions.

\subsection{The Heisenberg Drinfeld double}

The Heisenberg double of the positive and negative Mukai halves is
\begin{equation}
\label{eq:drinfeld-double-k3}
  D(\Yplus)=\Yplus\otimes\Yminus
\end{equation}
with the Mukai pairing identifying the two halves. This is the full
abelian K3 Yangian of Definition \ref{def:abelian-k3-yangian}. The compact
$K3\times E$ Hall-Drinfeld double belongs to the BKM branch and requires
extra compact Hall data; see
\texttt{chapters/examples/k3\_chiral\_algebra.tex},
Remark \texttt{rem:k3ca-hall-drinfeld-scope}.

\subsection{Stable envelopes at ADE/Kummer charts}

Maulik-Okounkov stable envelopes require an equivariant symplectic
resolution with a torus action. A generic K3 surface has
$\Aut(S)^\circ=0$, so this equivariant input is absent globally. At ADE
or Kummer charts the local surface model
$\widetilde{\C^2/\Gamma}$ carries the torus inherited from $\C^2$, and
the Hilbert scheme or Nakajima variety chart has stable envelopes.

\begin{proposition}[Local stable-envelope $R$-matrix]
\label{prop:mo-at-ade}
On an ADE/Kummer chart, the Maulik-Okounkov stable envelope produces a
rational $R$-matrix on the fixed-point basis of the local Nakajima
variety. Its Cartan restriction to Mukai weight spaces is diagonal and
matches \eqref{eq:r-matrix-diagonal} after identifying the chart weights
with the corresponding $h_i$.
\end{proposition}

\begin{proof}
Maulik-Okounkov's construction associates chamber transition matrices to
equivariant symplectic resolutions. For the ADE chart the fixed-point
basis is indexed by the quiver data, and the Cartan restriction is a
diagonal scalar action on weight spaces. The McKay identification embeds
those weights into $\Lmuk$. This proves chartwise agreement. A global K3
stable envelope would require a rank-$2$ torus action across the compact
surface, an input absent at generic K3 moduli.
\end{proof}

\begin{proposition}[Local compatibility of the three $R$-matrices]
\label{prop:three-routes}
On the common ADE/Kummer domain, the following constructions have the
same Cartan diagonal restriction:
\[
  \text{positive RTT half},\qquad
  \text{quiver CoHA shuffle half},\qquad
  \text{MO stable-envelope }R\text{-matrix}.
\]
The common restriction is the Mukai scalar family
$r_i(u)=(u-h_i)/(u+h_i)$.
\end{proposition}

\begin{proof}
The RTT half gives the scalar family by definition. The quiver shuffle
presentation inserts the same rational bond factor in the product. The
MO stable envelope records the same chamber wall transition on the
Cartan fixed-point basis. The three statements live on the local chart
where quiver, shuffle, and stable-envelope data are simultaneously
defined.
\end{proof}

\section{Holomorphic Chern-Simons boundary avatar}
\label{sec:hcs}

Let $C$ be a complex curve and consider holomorphic Chern-Simons theory
on $S\times C$ with action
\begin{equation}
\label{eq:hcs-action}
  S_{\mathrm{hCS}}[A]
  =
  \frac{1}{2\pi\hbar}
  \int_{S\times C}
  \Omega_S\wedge dz\wedge
  \Tr\!\left(A\,\bar\partial A+\frac{2}{3}A^3\right).
\end{equation}
Here $\Omega_S$ is the holomorphic symplectic form on K3 and $A$ is a
partial $(0,1)$-connection. Costello-Gwilliam factorisation supplies the
formal language for observables; the specific boundary identification
below is an additional conjectural quantisation statement.

\begin{conjecture}[hCS boundary algebra]
\label{conj:k3-yangian-as-boundary}
For a punctured disk $D^\times\subset C$, the boundary observable algebra
at the puncture is equivalent, as an $E_1$ chiral algebra, to the
abelian K3 Yangian:
\[
  \Obs_{\mathrm{hCS}}^\partial(S\times D^\times)
  \simeq
  \YgK.
\]
\end{conjecture}

\begin{proof}[Evidence]
The fibre contribution is the Mukai lattice
$H^0(S)\oplus H^2(S)\oplus H^4(S)$, and the disk contribution is the
one-dimensional chiral Heisenberg algebra. Their tensor product gives
the rank-$24$ current algebra \eqref{eq:k3-heisenberg-ope}. The rational
factor $(u-h_i)/(u+h_i)$ is the expected boundary propagator in the
$i$-th Mukai direction. Turning this evidence into a theorem requires
constructing the boundary condition, renormalised propagator, and
factorisation equivalence for hCS on $S\times C$.
\end{proof}

\begin{proposition}[Conditional OPE reading of stable envelopes]
\label{prop:mo-as-ope}
Assume Conjecture \ref{conj:k3-yangian-as-boundary}. Then the local
MO stable-envelope matrix on an ADE/Kummer chart is the OPE coefficient
for moving two boundary observables on disjoint punctured disks past one
another:
\[
  \Obs^\partial(S\times D_1^\times)\otimes
  \Obs^\partial(S\times D_2^\times)
  \xrightarrow{\ R_{\mathrm{MO}}(z_1-z_2)\ }
  \Obs^\partial(S\times(D_1^\times\cup D_2^\times)).
\]
\end{proposition}

\begin{proof}
Factorisation along disjoint disks gives a product of boundary
observables. Under the conjectural identification with $\YgK$, the
braiding is governed by the same diagonal rational $R$-matrix as in
Proposition \ref{prop:three-routes}.
\end{proof}

\subsection{Native $E_2$ object and $E_1$ mode algebra}

The native CY$_2$ output attached to $D^b\Coh(K3)$ is an $E_2$
factorisation algebra. The abelian K3 Yangian is the stage-two
one-dimensional specialisation, equivalently the mode algebra on a
formal disk, and is therefore $E_1$ as an associative chiral algebra. At
$d=3$, for $K3\times E$, the ambient algebra is $E_1$; the $E_2$ braiding
lives on the Drinfeld centre of its representation category.

\section{Branch separation and $K3\times E$ constants}
\label{sec:branch-separation}

\subsection{Kummer BFN chart}

At the Kummer point $T^4/\Z_2$ there are $16$ local $A_1$ charts. The
Braverman-Finkelberg-Nakajima Coulomb-branch construction applies to the
corresponding affine $A_1$ quiver variety and gives a local convolution
algebra. The comparison
\[
  A_{\mathrm{BFN}}(T^4/\Z_2,n)
  \rightsquigarrow
  \YgK|_{\mathrm{charge}\ n}
\]
is chartwise and conditional on deformation invariance through the
resolution of the orbifold singularities.

\subsection{Hall-Drinfeld and BKM branch}

The BKM superalgebra $\gg_{\Delta_5}$ attached to $K3\times E$ is the
Borcherds denominator object governed by Gritsenko's $\Delta_5$ at
$N=1$. Its compact quantum-group target is a Hall-Drinfeld/Borcherds
double built from a compact critical CoHA, a negative half, a Hopf
pairing, centre control, completion, and associator. The abelian K3
Yangian $\YgK$ is the Mukai-Heisenberg self-mirror branch. The two
branches share rank-$24$ Mukai data and have distinct algebraic
presentations.

\subsection{Four subscripted invariants on $K3\times E$}

Let $Y=K3\times E$. The four invariants used here are
\[
\begin{array}{c|c|l}
\text{invariant} & \text{value} & \text{construction}\\ \hline
\kappa_{\mathrm{cat}}(Y) & 0 &
\chi(\mathcal O_Y)=\chi(\mathcal O_{K3})\chi(\mathcal O_E)=2\cdot0\\
\kappa_{\mathrm{ch}}^{\mathrm{Heis}}(Y) & 3 &
\text{Heisenberg-Mukai chiral specialisation}\\
\kappa_{\mathrm{BKM}}(\Delta_5) & 5 &
c_1(0)/2=10/2,\ \text{Borcherds weight of }\Delta_5\\
\kappa_{\mathrm{fiber}}(Y) & 24 &
\operatorname{rk}\widetilde H(K3,\Z)=1+22+1
\end{array}
\]
Thus the $K3\times E$ spectrum is $\{0,3,5,24\}$. The value
$\chi(\mathcal O_{K3})=2$ is a fibre Euler characteristic, while the
total-space categorical invariant is $0$ by Kunneth. This matches the
local anchors
\texttt{chapters/examples/cy\_d\_kappa\_stratification.tex:135} and
\texttt{chapters/examples/k3e\_bkm\_chapter.tex:14290}.

\section{Proof boundaries}
\label{sec:proof-boundaries}

\begin{enumerate}
\item The diagonal RTT equation proves the positive current half. The
full Heisenberg central extension is the double pairing between positive
and negative modes.

\item The cited Schiffmann-Vasserot theorem
\texttt{arXiv:1705.07488} is a quiver CoHA generator theorem. It supplies
ADE/Kummer chart input. A compact global K3 CoHA with the stated double
requires extra gluing and orientation data.

\item The MO stable envelope is local to ADE/Kummer charts in this note.
Generic K3 moduli have $\Aut(S)^\circ=0$, so the necessary rank-$2$
equivariant torus is absent.

\item The coefficient
\[
  324=\dim\Sym^2(\Lmuk)+\dim\Lmuk=\binom{25}{2}+24
\]
is the $q^2$ Fock coefficient of the rank-$24$ Heisenberg module, also
the Hilbert-scheme oscillator count at charge $2$. A representation-count
interpretation at a root of unity would require a constructed small
quantum group quotient or braided finite category; this note constructs
only the generic rational Yangian object.

\item Root-of-unity specialisation is outside this construction. The rational parameter
$\hbar$ and the spectral parameters $h_i$ remain generic, subject only to
$\sum_i h_i=0$.
\end{enumerate}

\section{Open questions}
\label{sec:open-questions}

\begin{enumerate}
\item \emph{ADE nonabelian enhancement.}
At ADE walls the local shifted Yangian of the affine ADE quiver acts on
the corresponding Nakajima variety. The global assembly across K3 moduli
requires a deformation argument for the local $R$-matrices and their
Mukai embeddings.

\item \emph{Mukai-form super-extension.}
The Mukai form is symmetric on the full rank-$24$ lattice. The natural
envelope is the orthogonal Yangian $Y_\hbar(\mathfrak{so}(4,20))$ and
the programme-specific Hodge-parity super-extension
$Y_\hbar(\mathfrak{so}(4\mid20))$, whose odd part also carries the
symmetric Mukai form. The standard Kac superalgebra with a symplectic odd
form preserves a different bilinear datum from the Mukai lattice.

\item \emph{Compact Hall-Drinfeld double.}
The compact $K3\times E$ Hall-Drinfeld double requires the compact CoHA,
negative half, Hopf pairing, coproduct, completion, and associator. The
Borcherds denominator $\Delta_5$ gives the automorphic target, while the
compact Hall source remains an open construction.

\item \emph{Global stable-envelope replacement.}
A fibrewise construction over the K3 period domain would need to replace
the missing global torus by lattice-polarised wall crossing and
$O(\Lmuk)$-equivariance.

\item \emph{Mathieu twining.}
The $M_{24}$ action on K3 elliptic-genus data should act on the
Mukai-Heisenberg Fock space and commute with the Heisenberg current
algebra. Compatibility with the Frenkel-Kac cocycle is the first
vertex-level obstruction.
\end{enumerate}

\end{document}
